% !TEX program = xelatex
\documentclass[UTF8,fontset=fandol,a4paper,11pt]{ctexart}
\usepackage{../../_manual_common/hysim_manual}
\title{\textbf{HySim-XJTU-HRPES}\\[0.4em]{\Large model 模块实现契约手册}\\
{\large 数据语义、单位换算、有效容量与参数选择}}
\author{西安交通大学 HRPES 团队}
\date{\today}
\begin{document}
\maketitle
\begin{abstract}
本文档对应 \srcpath{src/model/} 与 \srcpath{include/hacdcpf/model/}，只记录代码实际执行的
数据规则和代数变换，不把元件容器描述扩展成潮流、优化或动态方程。规范投影实现虽然位于
\srcpath{src/model/}，其公共接口属于 \srcpath{include/hacdcpf/projection/}；因此仍归本模块说明，
但必须保留 rich、canonical 和结果恢复三类索引空间的区别。
\end{abstract}
\tableofcontents
\newpage

\section{模块定位、输入与输出}
输入是 \texttt{HybridPowerSystem} 及其 AC、DC、换流器和扩展组件容器；输出包括规范化系统、映射、
单位换算量、有效功率/容量和参数解析结果。本模块不求解电气状态。稳定对外 ID 是组件
\fld{index}；vector position 与 graph index 不是可交换的 ID。

\section{源码入口与数据结构}
\begin{longtable}{p{4.0cm}p{4.2cm}p{6.0cm}}
\toprule
源码 & 符号 & 实现职责\\\midrule
\srcpath{hybrid\_power\_system.hpp} & \texttt{HybridPowerSystem} & rich 工程语义容器\\
\srcpath{effective\_capacity.hpp} & \texttt{effective\_*} & 在役、缩放、可控性和容量回退\\
\srcpath{gfm\_norton\_contract.hpp} & \texttt{resolve\_gfm\_norton\_parameters} & GFM 参数优先级\\
\srcpath{unit\_conversion.cpp} & 单位换算函数 & 标幺、物理量和角度换算\\
\srcpath{project\_to\_canonical.hpp} & 投影与映射 & rich 到 canonical 及结果归因\\
\srcpath{model\_semantics.cpp} & 语义目录 & 组件在各工作流中的保真级别声明\\
\bottomrule
\end{longtable}

\section{与代码逐式一致的代数规则}
\subsection{缩放清洗和有效功率}
\srcpath{include/hacdcpf/model/effective\_capacity.hpp:sanitize\_scaling} 的
\srcpath{sanitize\_scaling} 逐字等价于
\begin{equation}
s_{eff}=\begin{cases}s,&\operatorname{isfinite}(s)\land s>0,\\0,&\text{其他情况。}\end{cases}
\end{equation}
AC/DC 负荷和静态发电的返回值分别来自同文件 :effective\_load\_p\_mw、:effective\_p\_mw（StaticGenerator）、:effective\_p\_mw（StaticGeneratorDC）：
\begin{align}
P_{load}^{eff}&=\begin{cases}P_{load}s_{eff},&\fld{in\_service}=true,\\0,&\text{否则，}\end{cases}\\
P_{sg}^{eff}&=\begin{cases}P_{sg}s_{eff},&\fld{in\_service}=true,\\0,&\text{否则，}\end{cases}\\
P_{dcsg}^{eff}&=\begin{cases}P_{set}s_{eff},&\fld{in\_service}=true,\\0,&\text{否则。}\end{cases}
\end{align}
这里不对 $P$ 本身取绝对值或截断；该符号行为与 helper 返回语句一致。

\subsection{有效容量回退}
静态 AC 发电的实现分支位于
\srcpath{include/hacdcpf/model/effective\_capacity.hpp:effective\_capacity\_mw}：
\begin{equation}
C_{sg}^{eff}=\begin{cases}
0,&\neg in\_service\lor s_{eff}=0,\\
\max(0,P s_{eff}),&\neg controllable,\\
\max(0,P^{cap}s_{eff}),&controllable,
\end{cases}
\end{equation}
其中 $P^{cap}=P_{max}$（若 $P_{max}>0$），否则用 $P_{rated}$（若 $P_{rated}>0$），再否则用 $P$。
DC 静态发电同式，但最后一级只在 $P_{max}$ 与 $P_{set}$ 间选择（:effective\_capacity\_mw）。
同步发电机容量（:effective\_capacity\_mw）为
\begin{equation}
C_g=\begin{cases}0,&\neg in\_service,\\P_{max},&P_{max}>0,\\\max(0,P_g),&\text{否则，}\end{cases}
\end{equation}
储能容量（:effective\_capacity\_mw）为
\begin{equation}
C_s=\begin{cases}0,&\neg in\_service,\\P_{max},&P_{max}>0,\\\max(0,P_{rated}),&\text{否则。}\end{cases}
\end{equation}

\subsection{GFM Norton 参数选择}
\srcpath{resolve\_gfm\_norton\_parameters}
（\srcpath{include/hacdcpf/model/gfm\_norton\_contract.hpp:resolve\_gfm\_norton\_parameters}）的选择规则为
\begin{align}
E&=\begin{cases}E_{gfm},&E_{gfm}>0,\\V_{ac,set},&\text{否则，}\end{cases}\\
\theta&=\frac{\pi}{180}\begin{cases}
\theta_{gfm},&E_{gfm}>0\lor\theta_{gfm}\ne0,\\
\theta_{ac,set},&\text{否则，}
\end{cases}\\
R_v&=\begin{cases}R_{gfm},&R_{gfm}\ne0,\\R_{conv},&\text{否则，}\end{cases}&
X_v&=\begin{cases}X_{gfm},&X_{gfm}\ne0,\\X_{sc},&\text{否则。}\end{cases}
\end{align}
函数随后返回 $\underline E=E e^{j\theta}$ 与 $Z_v=R_v+jX_v$。零值在这里具有“使用回退字段”的精确
语义，不能改写成非负约束或默认物理参数。

\section{规范投影与映射不变量}
投影可形式化为代码映射
\begin{equation}
(S_c,M)=\Pi(S_r),\qquad
M=(M_{bus}^{ac},M_{bus}^{dc},M_{branch},M_{component}).
\end{equation}
该式只表示函数输入输出，不声称所有组件均为无损等价。允许结果恢复的最小不变量是：
\begin{align}
M_{bus}^{ac}&:\text{canonical AC position}\to\text{原始 AC 稳定 ID},\\
M_{bus}^{dc}&:\text{canonical DC position}\to\text{原始 DC 稳定 ID},\\
M_{branch}&:\text{canonical branch}\to\texttt{BranchRef}.
\end{align}
AC 与 DC 可拥有相同整数 ID，因此映射键必须包含 domain。聚合、零阻抗合并或三绕组展开时不保证
组件数量守恒；结果只能经对应映射恢复，不能按 vector position 回填。

\section{边界条件、复杂度与异常处理}
容器扫描和 helper 为 $O(n)$ 或单元素 $O(1)$；投影成本至少与组件数线性，并受图合并与映射构建影响。
非有限/负缩放被精确清零；缺失 GFM 专用字段按上述零值规则回退；无效电气引用不由本模块 helper 自动
修复，应交给 \texttt{validation}。投影生成的近似组件必须通过映射或限制字段保留来源，不能宣称 rich
模型全部精确覆盖。

\section{接口、测试与工业评价}
上游是 I/O 和人工建模；下游包括 validation、graph、PF、OPF、动态、可靠性和市场。测试应覆盖：
非有限缩放、零缩放、可控/不可控容量分支、GFM 每个零值回退分支、AC/DC 同号 bus ID、投影往返、
零阻抗合并和结果归因。评价指标为映射完备率、往返字段保真率、未知组件拒绝率、单位一致性和跨模块
有效容量一致性。

\section{适用范围、局限与改进方向}
本模块提供语义和转换，不提供物理可解性保证。历史
\srcpath{component\_models\_math\_audit.tex} 是问题审计，不是当前行为规范。后续改进应把每一类 rich
组件的 projection fidelity 与结果恢复证书变成机器可读表，并由测试自动检查。
\section{跨模块工业级技术评价基线}

本章规定本手册所述模块进入工程算例、批量研究或生产服务前必须共同满足的评价基线。
具体模型公式、变量、约束和专用算法以正文相应章节为准；本章不扩大源码已经实现的范围。

\subsection{输入、输出、边界条件与数据结构审计}
输入审计应逐项检查有限性、单位、上下界、枚举值和引用完整性。系统对象必须先按所需层级完成
校验；AC 与 DC 母线采用域限定映射，组件稳定 \texttt{index}、向量位置和临时图索引不得混用。
输出审计必须记录单位、索引空间、模型范围、收敛或可行状态、后端、容差、迭代/节点计数和告警。
空系统、空时域、孤岛、重复 ID、零容量、退化约束与不可用可选后端均属于边界测试集。

\subsection{工程假设与数值稳定性分析}
模型使用的平衡、线性化、稳态、独立性、概率分布、时间离散或网络等值假设必须与应用场景匹配。
数值评价至少包含变量缩放、绝对/相对残差、可行性违反、条件数或枢轴质量、步长/阻尼、迭代停滞
和随机种子复现性。若采用正则化、平滑、回退、松弛或近似，应同时报告触发条件、参数值、对主指标
的敏感性以及是否改变模型口径；不得用“已返回结果”替代收敛或有效性证明。

\subsection{复杂度与资源分析}
令 $n_x$、$n_c$、$n_t$、$n_s$、$|V|$、$|E|$ 分别表示变量、约束、时段、场景、节点和边数。
报告应分别给出建模、求解、后处理的时间与内存。图遍历通常为 $O(|V|+|E|)$；逐时段/逐场景
流程至少线性增长为 $O(n_t n_s)$ 次子问题调用；稠密线性代数最坏可达 $O(n_x^3)$，稀疏方法则应
报告结构非零数、填充率和因子分解峰值内存；MILP/NLP 不给出虚假的多项式最坏界，而报告节点数、
间隙、时限和实测扩展曲线。复杂度结论必须基于固定硬件、固定后端和可复现算例族。

\subsection{异常工况处理}
非法输入应在进入求解器前返回结构化错误；不可行、无界、不收敛、时限、内存不足和后端缺失必须
相互区分。部分结果只有在对应 \fld{ValidityFlags}、\fld{model\_scope}、
\fld{model\_limitations} 或等价字段允许时才可使用。异常路径不得静默丢弃 AC/DC 资产、制造平衡
节点、把 NaN 改写为零或把近似解标为全局最优；系统原始对象应保持可恢复和可重试。

\subsection{与其他模块的接口关系}
接口链路遵循“工程语义模型 $\to$ 校验 $\to$ 投影/装配 $\to$ 求解或分析 $\to$ 结果回投影”。
跨模块传递时保留稳定 ID、单位、符号、时间基准和场景身份；消费潮流、OPF、动态或可靠性结果前，
先检查其收敛、覆盖范围和有效性标志。HTTP/JSON/Python 层只做语义保持适配，不重新解释求解结果。

\subsection{测试与验证建议}
最低验证矩阵包括：手算小例、标准算例、边界/异常输入、随机性质测试、算法或后端交叉校核、
序列化往返、稳定 ID 回投影、AC/DC 同号母线、重复运行确定性和规模扩展测试。数值算法还应加入
残差独立重算、雅可比有限差分或自动微分校核；近似模型应与高保真模型比较并给出误差分布，而非
只报告单个平均值。回归报告必须列明构建配置、算例版本、容差、通过/失败/跳过数及未验证范围。

\subsection{工业级评估指标}
\begin{itemize}
  \item \textbf{正确性}：守恒残差、约束最大违反、解析/基准误差、结果回投影一致性；
  \item \textbf{鲁棒性}：收敛率、不可行识别率、异常分类准确率、扰动敏感性与随机复现率；
  \item \textbf{性能}：端到端时延、吞吐量、峰值内存、并行效率与规模增长斜率；
  \item \textbf{可追溯性}：输入模式版本、稳定 ID 覆盖率、模型范围/限制字段完整率、告警可定位率；
  \item \textbf{工程有效性}：与标准、外部引擎或现场数据的偏差，以及误判/漏判对业务指标的影响。
\end{itemize}
指标应预先给出验收阈值和失败处置；未达到阈值时只能标记为实验性或受限适用。

\subsection{适用范围、局限性与后续改进}
本手册只评价当前源码覆盖的模型、后端和数据结构，不对未建模物理过程、未安装求解器或未执行的
外部验证作保证。后续改进应按“缺口 $\to$ 可量化指标 $\to$ 源码/测试变更 $\to$ 独立验证”闭环，
优先消除会造成错误结论的语义缺口，其次提升数值鲁棒性与性能，最后扩展模型范围。

\end{document}
